\subsection{Positive coefficient tails have no higher graded dual obstruction}
\label{arch136:tail-ext}

\paragraph{Archive scope and conventions.}
This is a coefficient calculation for the coherent Moore construction
at cote136, physical p.~117/footer116/scan117 and physical
p.~121/footer120/scan121.  It applies to the coefficient complex in
ARCH-136-DERIVED-TAIL-COMPARISON.  All arguments are proved below;
there is no imported theorem about the homological dimension of a
nonnoetherian divided-power algebra.  Write
\[
 R=\mathbb Z\quad\hbox{or}\quad R=\mathbb Z_{(p)},\qquad
 S=\Gamma_R(T)=\bigoplus_{w\geq0}R T^{[w]},\qquad
 I_M=\bigoplus_{w\geq M}R T^{[w]},\quad M\geq1.
                                                        \tag{TE.1}
\]
Weights, rather than form degrees, grade these modules.  The notation
\(M\langle v\rangle\) increases every weight by \(v\).
Let \(\mathsf{Gr}_{\geq0}(S)\) denote the abelian category of
\(S\)-modules concentrated in nonnegative weights, with maps of
weight zero.  Its projective generators are
\(S\langle v\rangle\), \(v\geq0\).
Every assertion of injectivity below is in this category.

\paragraph{ARCH-136-GRADED-INJECTIVES.}
The exact sequence
\[
 0\longrightarrow S\longrightarrow E=\mathbb Q[T]
       \longrightarrow J=\mathbb Q[T]/S\longrightarrow0
                                                        \tag{TE.2}
\]
is an injective resolution of length one in
\(\mathsf{Gr}_{\geq0}(S)\).  The embedding uses
\(T^{[w]}=T^w/w!\).  This statement holds both over
\(\mathbb Z\) and over \(\mathbb Z_{(p)}\).

\begin{proof}
We first spell out the version of the ideal criterion used here.
A nonnegatively graded module \(V\) is injective if every
weight-zero map
\[
                 L\langle v\rangle\longrightarrow V,
       \qquad L\subset S\text{ a homogeneous ideal},\quad v\geq0,
                                                        \tag{TE.3}
\]
extends to \(S\langle v\rangle\).  Indeed, extend a map on a
submodule as far as possible by the usual maximal-extension argument.
If a homogeneous element \(x\) of weight \(v\geq0\) remains outside
the domain, the elements \(s\) for which \(sx\) already lies in the
domain form a homogeneous ideal \(L\).  An extension in (TE.3)
assigns a compatible image to \(x\), enlarging the domain, a
contradiction.  This proves the criterion directly in the stated
category.

For \(E\), a map in (TE.3) from a nonzero ideal has the form
multiplication by a rational multiple of \(T^v\).  To see this,
choose one nonzero homogeneous \(a\in L\).  For every homogeneous
\(b\in L\), linearity gives \(a f(b)=b f(a)\), so
\(f(b)=b f(a)/a\) in \(\mathbb Q(T)\).  Its weight makes
\(f(a)/a\in\mathbb Q T^v\subset E\), since \(v\geq0\).
This multiplier extends the map to \(S\langle v\rangle\).
The zero ideal causes no difficulty.  Thus \(E\) is injective.

Now take \(R=\mathbb Z_{(p)}\).  For
\(D=p^r-1\), put \(A_D=S/I_{D+1}\), and truncate any
nonnegative graded module by quotienting its weights above \(D\).
The coefficient of \(T^{[D]}\) defines a perfect pairing on
\(A_D\).  In its divided-power basis the only nonzero entries
of that pairing are the anti-diagonal entries
\[
                         \binom Da\in R^\times,
                     \qquad 0\leq a\leq D.              \tag{TE.4}
\]
For example their nonvanishing modulo \(p\) follows by expanding
\((1+X)^{p^r-1}=(1+X^{p^r})/(1+X)\) in characteristic \(p\).
Consequently there is a graded \(A_D\)-linear isomorphism
\[
 J_{\leq D}=A_D\otimes_R(\mathbb Q/R)
   \simeq\operatorname{Hom}_R(A_D,\mathbb Q/R)\langle D\rangle.
                                                        \tag{TE.5}
\]
Here the graded dual on the right has weights from \(-D\) to zero
before its displayed shift.  The module \(\mathbb Q/R\) is injective
over \(R\): for any nonzero ideal \(p^aR\), extending a map from
that ideal just asks to divide one element by \(p^a\), which is
possible in \(\mathbb Q/R\).  The same maximal-extension proof
applies.  The coinduction adjunction therefore makes (TE.5)
injective as a graded \(A_D\)-module.  In particular, taking
weight-zero Hom into it is exact.

Let \(f:L\langle v\rangle\to J\), with \(v\geq0\).
For each \(D=p^r-1\geq v\), truncation gives an inclusion of
\(A_D\)-modules
\[
 (L\langle v\rangle)_{\leq D}
      \subset (S\langle v\rangle)_{\leq D}.             \tag{TE.6}
\]
These source modules truncate at coefficient weight \(D-v\);
they are nevertheless \(A_D\)-modules because their weights are
nonnegative.  By (TE.5), the truncated map extends across (TE.6).
Such an extension is determined by the image
\(c\in J^v\) of the free generator.  Let \(X_D\subset J^v\)
be the set of these possible values.  These sets are nonempty
and decrease as \(D\) increases through the stated sequence.

If \(L\ne0\), fix a nonzero homogeneous element
\(a=\alpha T^{[w]}\in L\), \(\alpha\in R\setminus\{0\}\).
Identify \(J^v\) and \(J^{w+v}\) with \(\mathbb Q/R\)
using their divided-power bases.  Once \(D\geq w+v\), every
\(c\in X_D\) satisfies the single equation
\[
             \alpha\binom{w+v}{v}\,c=f(a).              \tag{TE.7}
\]
Multiplication by the nonzero scalar on the left has finite kernel
on \(\mathbb Q/R\).  Thus all these nonempty nested sets lie
in one finite coset.  Their intersection is nonempty.  A value
in that intersection defines an extension on
\(S\langle v\rangle\), since every element of \(L\) appears
in some finite truncation.  The ideal criterion proves that
\(J\) is injective over \(\Gamma_{\mathbb Z_{(p)}}(T)\).

For \(R=\mathbb Z\), decompose the graded torsion module as
\[
 J=\bigoplus_p J_{(p)},\qquad
 J_{(p)}\simeq\mathbb Q[T]/\Gamma_{\mathbb Z_{(p)}}(T).
                                                        \tag{TE.8}
\]
Project a map \(f\) in (TE.3) to each summand.  It extends
uniquely to a map on \(L\otimes\mathbb Z_{(p)}\), since
integers prime to \(p\) act invertibly on the target.  The local
argument supplies an extension value \(c_p\in J_{(p)}^v\).
For the fixed nonzero element \(a=\alpha T^{[w]}\), now with
\(\alpha\in\mathbb Z\setminus\{0\}\), equation (TE.7)
shows that \(c_p=0\) outside the finite set of primes dividing
\(\alpha\binom{w+v}{v}\) or occurring in the order of \(f(a)\).
For those other primes multiplication by that integer is invertible
and \(f(a)_p=0\).  Hence \(\sum_p c_p\) is an actual element
of \(J^v\), and its multiplication map extends all components of
\(f\) simultaneously.  This proves global injectivity of \(J\)
and completes (TE.2).
\end{proof}

\paragraph{ARCH-136-TAIL-EXT-VANISHING.}
For \(M\geq1\) and \(v\geq0\), restriction of multiplication
maps gives the canonical identifications
\[
 \operatorname{Hom}_{\mathrm{gr},S}(I_M\langle v\rangle,S)=S^v,
 \qquad
 \operatorname{Ext}^{i}_{\mathrm{gr},S}(I_M\langle v\rangle,S)=0
                         \quad(i>0).                    \tag{TE.9}
\]
Moreover
\[
 \operatorname{Ext}^{i}_{\mathrm{gr},S}
       ((S/I_M)\langle v\rangle,S)=0\qquad(i\geq0).      \tag{TE.10}
\]
Here Ext can equally be computed in all graded modules: each of
these nonnegative modules has a free resolution with generator
weights nonnegative, so the same resolution computes both versions.

\begin{proof}
Any map \(I_M\langle v\rangle\to S\), regarded as a map to
\(E\), is multiplication by \(cT^{[v]}\), \(c\in\mathbb Q\).
Its value on \(T^{[w]}\), \(w\geq M\), is
\[
                  c\binom{w+v}{v}T^{[w+v]}.             \tag{TE.11}
\]
For any relevant prime \(p\), choose \(r\) so large that
\(w=p^r-1-v\geq M\).  The binomial coefficient in (TE.11)
is a \(p\)-unit by (TE.4).  Integrality of the image therefore
forces \(c\in\mathbb Z_{(p)}\).  This proves \(c\in R\)
in the local case and, by all primes, in the integral case.
Conversely every such multiplier works.  This proves the Hom claim.

Use the injective resolution (TE.2).  Any map
\(I_M\langle v\rangle\to J\) extends to
\(S\langle v\rangle\), since \(J\) is injective.  Its one
extension value in \(J^v\) lifts to \(E^v\).  Thus
\(\operatorname{Hom}(I_M\langle v\rangle,E)\to
\operatorname{Hom}(I_M\langle v\rangle,J)\) is onto,
which gives vanishing of Ext in degree one.  Higher degrees vanish
by (TE.2).  Finally apply Hom to
\(0\to I_M\langle v\rangle\to S\langle v\rangle\to
(S/I_M)\langle v\rangle\to0\).  The degree-zero restriction
map is the isomorphism already proved, while a map from the quotient
to the domain \(S\) is zero.  This gives (TE.10).
\end{proof}

\paragraph{Canonical comparison and all finite moduli.}
Let \(P_\bullet\to I_M\) be any graded free resolution with
generator weights at least \(M\), and put
\[
 K^i=\operatorname{Hom}_{\mathrm{gr},S}(P_i\langle v\rangle,S).
\]
The inclusion of the tail into \(S\), composed with the resolution
augmentation, defines the canonical quasi-isomorphism
\[
                         S^v[0]\longrightarrow K.       \tag{TE.12}
\]
Each \(K^i\) is a product of weight components of \(S\), hence
is torsion-free over \(R\).  Reduction of that product modulo
any integer \(N\geq1\) is the product of its reductions: choose
coordinate lifts to see surjectivity, and its kernel is exactly
\(N K^i\).  Reduction of the free resolution is exact because
\(I_M\) and every syzygy are torsion-free over \(R\), hence
flat over this principal ideal domain.  Consequently
\[
 \begin{split}
 K/NK&=\operatorname{Hom}_{\mathrm{gr},S/NS}
                 (P_\bullet/NP_\bullet\langle v\rangle,S/NS),\\
 H^i(K/NK)&=
 \begin{cases}S^v/NS^v,&i=0,\\0,&i>0.
 \end{cases}
 \end{split}                                             \tag{TE.13}
\]
The last assertion also follows immediately by applying the cone
of multiplication by \(N\) to (TE.12).  With the cochain convention
\(P_N(K)^r=K^r\oplus K^{r+1}\),
\(d(x,y)=(dx-Ny,-dy)\), the map
\((x,y)\mapsto x\bmod N\) identifies this cone, up to
quasi-isomorphism, with \(K/NK\).  Thus it has only degree-zero
cohomology, canonically \(S^v/NS^v\).  These comparisons are
compatible with every map of coefficient tails and every finite
modulus; no inverse-limit argument or tensoring an infinite Hom
product with \(\mathbb Q\) has been used.

\paragraph{Exact consequence and remaining gate.}
The higher coefficient Ext groups in the derived tail construction
are zero for the nonnegative shifts that occur there, over both
\(\mathbb Z\) and \(\mathbb Z_{(p)}\).  This removes a proposed
linear coherence obstruction and permits graph comparisons to use
ordinary coefficient groups.  It does not assert that a subsequent
free PD realization preserves vertical exactness.  In particular,
it is compatible with ARCH-136-ACTUAL-TAIL-SYZYGY: a coefficient
resolution can have vanishing dual Ext while its bottom parameter
module has an even-PD disk with nonzero degree-zero torsion.
Negative shifts, arbitrary coefficient target modules, and arbitrary
rings are outside the theorem.

\paragraph{Related primary source and precise distinction.}
Nagpal--Snowden, \emph{The module theory of divided power algebras},
arXiv:1606.03431v2 (18 February 2018), Theorem~7.1(a), proves
characteristic-\(p\) graded self-injectivity.  Their Theorems~7.9
and~7.11 impose finite-presentation hypotheses for the stated stronger
Ext bounds.  Remark~7.12 warns that a negative shift obstructs an
unrestricted injectivity claim; Theorem~8.3 and Remark~8.4 record the
completion issue for arbitrary finitely presented targets.  The relevant
v2 printed pages are 39--44 and 45--46.  The special integral,
nonnegative-category formulation above has its own complete proof;
no novelty claim or unrestricted extension of those results is made.

\paragraph{Verification status.}
The author, the parent, and the independent peer on the prime-degree
branch have each read this complete saved proof and reported PASS.
Their checks include the exact nonnegative graded category, both
injectives, shifted finite truncations, finite extension cosets, global
finite-prime support, all tail Hom and Ext claims, and every finite
modulus cone.  The separate check events record their actual scope.
No main connectedness conclusion is claimed.
